% Supplementary material draft — REVISION RESERVE (user 2026-11-13: not submitted with the initial manuscript;
% kept ready in case reviewers ask for more detail). Tables are generated by make_supp_tables.py into tables_generated.tex.
% Build: bash build_supp.sh  (pdflatex + bibtex, output supp_main.pdf)
\documentclass[10pt]{article}
\usepackage[T1]{fontenc}
\usepackage[scaled=0.92]{helvet}
\usepackage{mathptmx}
\usepackage[a4paper,margin=2.2cm]{geometry}
\usepackage{booktabs,graphicx,xcolor,amsmath,caption}
\usepackage[hidelinks]{hyperref}
\captionsetup{font=small,labelfont=bf}
\renewcommand{\thetable}{S\arabic{table}}
\renewcommand{\thefigure}{S\arabic{figure}}
\newcommand{\todo}[1]{\textcolor{red}{[TODO: #1]}}

\title{Supplementary Material\\[4pt]\large Clean-Teacher Distillation on Corrupted Sensor Inputs for Robust Multimodal Semantic Segmentation}
\author{Jemo Maeng and Kyoobin Lee}
\date{}

\begin{document}
\maketitle

This document complements the main paper with tables that are too detailed for the main text. Table and figure numbers without the prefix S refer to the main paper. All numbers of our model are the mean over the training runs stated in the main paper unless a caption says otherwise.

\section{Missing sensors per subset on DELIVER}
Table~\ref{tab:s_emm} expands the entire-missing average of Table~5 (benchmark protocol) into the 15 subsets of removed sensors, measured with the public benchmark code on the validation set. The released CMNeXt weights reproduce the benchmark's reported values, and the averages in the last rows are those of Table~5. CMNeXt collapses whenever depth is removed (22.9 with depth alone missing, 3.2 or less with depth and RGB missing), whereas our model keeps at least 24.9 in every subset and matches or exceeds its clean score when event, LiDAR, or both are removed (71.5 to 72.0 against 71.3).

\section{MUSES test results per condition and class}
Table~\ref{tab:s_muses_cond} gives the official test scores of our three MUSES models per condition, next to the per-condition results that MM SAM-adapter reports for the same two-sensor settings. Table~\ref{tab:s_muses_cls} gives the per-class scores of the same submissions. Each of our entries is one submission (the run with the higher validation score), so differences of about one point between models or conditions are within what a different training run could change.

\section{MCubeS per class}
Table~\ref{tab:s_mcubes_cls} compares, per material class, our full model with the clean-trained model without the high-resolution path that serves as its teacher.

\section{Dependence on depth per condition}
Table~\ref{tab:s_depth} lists the per-condition drops plotted in Fig.~7 of the main paper.

\section{Planned for revision: C2D2 on other fusion architectures}
\todo{Only if reviewers ask (lab-plan DRN-261029-02/03, revision reserve). Apply the full C2D2 scheme (clean pass, degraded pass, frozen clean teacher) to DGFusion and CAFuser and report clean mIoU, EMM average, and the per-condition depth drop next to their single-pass re-training of Table~\ref{tab:s_depth} and of the re-trained comparison in the main paper.}

% generated by make_supp_tables.py — do not edit by hand
\begin{table}[t]
\caption{Entire-missing evaluation on the DELIVER validation set per removed-sensor subset (mIoU, \%), public benchmark code~\cite{mmssbench2025}. R: RGB, D: depth, E: event, L: LiDAR.}
\label{tab:s_emm}
\centering\small
\begin{tabular}{lcc}
\toprule
Removed sensors & CMNeXt (released) & Ours (mean of 3 runs) \\
\midrule
none (clean) & 66.33 & \textbf{71.34} \\
R & 50.98 & \textbf{62.28} \\
D & 22.92 & \textbf{63.13} \\
E & 66.38 & \textbf{71.53} \\
L & 66.27 & \textbf{71.62} \\
R + D & 3.15 & \textbf{36.93} \\
R + E & 50.83 & \textbf{60.71} \\
R + L & 50.80 & \textbf{62.23} \\
D + E & 22.50 & \textbf{63.33} \\
D + L & 22.91 & \textbf{63.24} \\
E + L & 66.33 & \textbf{71.97} \\
R + D + E & 2.86 & \textbf{24.93} \\
R + D + L & 3.16 & \textbf{31.29} \\
R + E + L & 50.59 & \textbf{60.25} \\
D + E + L & 22.49 & \textbf{62.98} \\
\midrule
Average over 15 subsets & 37.90 & \textbf{58.52} \\
E(.2) & 54.46 & \textbf{67.15} \\
E(.1) & 60.41 & \textbf{69.46} \\
E(.05) & 63.38 & \textbf{70.45} \\
\bottomrule
\end{tabular}
\end{table}

\begin{table*}[t]
\caption{MUSES test set per condition (mIoU, \%, official evaluation server). MM SAM-adapter values are from~\cite{mmsamadapter2025} (Table~6); ours are one submission each.}
\label{tab:s_muses_cond}
\centering\small
\begin{tabular}{lccccccc}
\toprule
Model & All & Day & Night & Clear & Fog & Rain & Snow \\
\midrule
MM SAM-adapter, RGB-L & \textbf{81.07} & \underline{83.34} & 74.97 & 80.82 & \underline{74.12} & \textbf{80.00} & \textbf{80.85} \\
MM SAM-adapter, RGB-E & 79.92 & \textbf{83.39} & 72.38 & \textbf{81.52} & 68.97 & 78.92 & 78.72 \\
\midrule
Ours, RGB-L & \underline{80.46} & 81.53 & \textbf{77.56} & 80.55 & 74.00 & 79.15 & \underline{80.07} \\
Ours, RGB-E & 79.91 & 82.68 & 74.27 & \underline{80.90} & 72.52 & 78.60 & 77.91 \\
Ours, RGB-L-R-E & 80.01 & 81.97 & \underline{75.83} & 79.86 & \textbf{79.23} & \underline{79.34} & 78.84 \\
\bottomrule
\end{tabular}
\vspace{4pt}

\begin{tabular}{lcccccccc}
\toprule
Model & Clear day & Clear night & Fog day & Fog night & Rain day & Rain night & Snow day & Snow night \\
\midrule
Ours, RGB-L & 81.66 & \textbf{77.04} & 72.85 & \underline{67.84} & 78.86 & \textbf{74.36} & 72.06 & \textbf{78.16} \\
Ours, RGB-E & \textbf{83.15} & 74.45 & \underline{74.33} & 64.41 & \underline{79.83} & 72.65 & \textbf{72.72} & 72.00 \\
Ours, RGB-L-R-E & \underline{82.02} & \underline{74.81} & \textbf{80.03} & \textbf{72.54} & \textbf{80.38} & \underline{74.14} & \underline{72.69} & \underline{77.12} \\
\bottomrule
\end{tabular}
\end{table*}

\begin{table}[t]
\caption{MUSES test set per class (IoU, \%, official evaluation server, all conditions).}
\label{tab:s_muses_cls}
\centering\small
\begin{tabular}{lccc}
\toprule
Class & RGB-L & RGB-E & RGB-L-R-E \\
\midrule
Road & \textbf{97.23} & \underline{97.17} & 97.09 \\
Sidewalk & \textbf{87.26} & \underline{87.14} & 86.99 \\
Building & \textbf{93.59} & 93.01 & \underline{93.55} \\
Wall & \underline{82.04} & 78.73 & \textbf{82.42} \\
Fence & \textbf{67.90} & 64.12 & \underline{66.26} \\
Pole & 65.18 & \textbf{65.58} & \underline{65.34} \\
Traffic light & \textbf{71.93} & 71.56 & \underline{71.61} \\
Traffic sign & \underline{75.18} & 73.00 & \textbf{75.27} \\
Vegetation & \textbf{89.87} & 88.57 & \underline{89.70} \\
Terrain & \textbf{79.13} & 76.43 & \underline{77.74} \\
Sky & \underline{96.89} & 96.52 & \textbf{97.10} \\
Person & \underline{72.92} & \textbf{73.11} & 72.79 \\
Rider & 59.51 & \textbf{62.28} & \underline{60.05} \\
Car & \textbf{94.64} & \underline{94.38} & 94.35 \\
Truck & \textbf{77.64} & \underline{77.63} & 77.33 \\
Bus & 95.38 & \underline{95.70} & \textbf{95.77} \\
Train & \underline{93.96} & \textbf{94.43} & 93.12 \\
Motorcycle & \textbf{59.22} & \underline{58.81} & 56.74 \\
Bicycle & \underline{69.26} & \textbf{70.19} & 67.02 \\
\midrule
mIoU & \textbf{80.46} & 79.91 & \underline{80.01} \\
\bottomrule
\end{tabular}
\end{table}

\begin{table}[t]
\caption{MCubeS test set per class (IoU, \%, mean of two runs): the clean-trained model without the high-resolution path and our full model.}
\label{tab:s_mcubes_cls}
\centering\small
\begin{tabular}{lcc}
\toprule
Class & Clean-trained, no high-res. & Ours \\
\midrule
Asphalt & 88.81 & \textbf{89.62} \\
Concrete & \textbf{55.99} & 55.80 \\
Metal & 64.56 & \textbf{65.10} \\
Road marking & 74.08 & \textbf{77.13} \\
Fabric & \textbf{45.16} & 42.01 \\
Glass & 61.72 & \textbf{61.84} \\
Plaster & 0.19 & \textbf{1.51} \\
Plastic & 50.53 & \textbf{52.39} \\
Rubber & 27.18 & \textbf{32.89} \\
Sand & 61.97 & \textbf{64.91} \\
Gravel & 67.91 & \textbf{72.28} \\
Ceramic & 40.21 & \textbf{42.30} \\
Cobblestone & 73.47 & \textbf{76.41} \\
Brick & \textbf{52.65} & 51.27 \\
Grass & \textbf{62.72} & 60.90 \\
Wood & \textbf{53.03} & 50.52 \\
Leaf & \textbf{77.78} & 77.25 \\
Water & 85.30 & \textbf{86.31} \\
Human & 35.98 & \textbf{38.40} \\
Sky & 96.21 & \textbf{96.67} \\
\midrule
mIoU & 58.78 & \textbf{59.78} \\
\bottomrule
\end{tabular}
\end{table}

\begin{table*}[t]
\caption{Drop in mIoU (points) on the DELIVER test set when depth is removed (zero-filled), per condition; lower is better. Data of Fig.~7 of the main paper.}
\label{tab:s_depth}
\centering\small
\begin{tabular}{lcccccc}
\toprule
Model & Sun & Cloud & Rain & Fog & Night & Low exp. \\
\midrule
DGFusion (our reproduction) & 23.1 & 24.9 & 27.6 & 25.9 & 32.6 & 37.1 \\
Ours, clean training, no high-res. path (2 runs) & 14.6 & 14.8 & 14.3 & 14.4 & 17.9 & 25.0 \\
Ours (3 runs) & 1.3 & 1.4 & 1.5 & 0.6 & 4.3 & 9.6 \\
\bottomrule
\end{tabular}
\end{table*}


\bibliographystyle{IEEEtran}
\bibliography{../latex/references}
\end{document}
